\chapter{引论}

\section{选题的研究现状}

中国是一个人口多而入学率低的国家，每年大概有一半的高中生升入大学，但我国高中教授的导数内容与其他国家相比差异较大。

在导数教育方面，我国强调实用性和快捷性教育。为了防止学生忽视导数的数学思想价值和表达形式，教师应避免无用的形式化练习。 



\section{研究意义}

导数作为高中数学的一部分内容，在不同学习阶段的内涵和外延是不同的，不同教材对导数的概念教学及其导数的教学设置和教学要求是不同的。

学生应着眼于用导数的工具性及其思想方法来解决数学学习中的相关问题，即导数在研究曲线的切线方程和函数的若干性质以及最优化问题方面提供捷径，因此导数在高中数学阶段越来越凸现其重要性。

课本选修2中，函数的一些性质（如单调性、极值、最值等等）都是用求导法研究的。但是，在平常的解题中，学生们还是会下意识地选择最麻烦的方式来解决问题。本文的目的就在于使学生真正学会利用导数作为解题工具为高中数学大题服务，并锻炼学生的逻辑思维能力和解题能力，为未来接受高等数学教育作准备。